\documentclass[a4paper,12pt]{article}

%	Tous les packages pour francisation compl\' ete
%\usepackage[applemac]{inputenc}
\usepackage[T1]{fontenc}
\usepackage{lmodern,textcomp}
\usepackage[frenchb]{babel}
\usepackage{multicol}
\usepackage{amsfonts,amsmath,amssymb,mathrsfs}
\usepackage{geometry}
\usepackage{aeguill}
\usepackage{hyperref}
\usepackage{array}
\usepackage{multirow}
\usepackage{asymptote}
\usepackage{graphicx}
\usepackage{tabularx}
\usepackage{enumerate,tablists, enumitem}
\usepackage{eurosym}
\usepackage{pst-solides3d}
\usepackage{multirow}
\usepackage{extsizes}
\usepackage{calc}
\usepackage{boites}
\usepackage{color,colortbl}
\usepackage{fancyhdr}
\usepackage{stmaryrd}
\usepackage{pifont}
%\usepackage{pst-all,pstricks,pstricks-add,pst-plot,pst-eucl,pst-text,pst-tree,pst-math,pst-eps}

%\usepackage{auto-pst-pdf}
\makeatletter
\let\Test@pr@shipout\pr@shipout%% save the original definition
\let\Test@shipout\shipout
\makeatother
\usepackage{tikz,tkz-tab}
\makeatletter
\AtBeginDocument{%
  \let\pr@shipout\Test@pr@shipout%% restore it 
  \let\shipout\Test@shipout
}
\makeatother

\newtheorem{definition}{D\'efinition}
\newtheorem{theoreme}{Th\'eor\`eme}

\newcommand{\C} {\ensuremath{\mathbb{C}}}
\newcommand{\R} {\ensuremath{\mathbb{R}}}
\newcommand{\Z} {\ensuremath{\mathbb{Z}}}
\newcommand{\N} {\ensuremath{\mathbb{N}}}
\newcommand{\ds} {\displaystyle}
\newcommand{\p}{\par}
\newcommand{\abs}[1]{\left |#1\right |}
\newcommand{\eq}[1]{\underset{#1}{\sim}}
\newcommand{\tend}[1]{\underset{#1}{\longrightarrow}}
\newcommand{\norme}[1]{\left\Vert #1\right\Vert}
\newcommand{\V}[1]{\overrightarrow{#1}} %vecteur

\definecolor{gris1}{gray}{0.85}
\definecolor{gris2}{gray}{0.65}

\setlength{\textwidth} {19cm}
\setlength{\textheight} {26cm}
\addtolength{\topmargin} {-2.5cm}
\setlength{\oddsidemargin} {-1.5cm}
\setlength{\evensidemargin} {-1.5cm}

\newcommand{\encad}[1]{\begin{center}%
\fcolorbox{black}{gray!20}{\begin{minipage}[t]{0.8\linewidth}%
#1\end{minipage}}\end{center}}


\pagestyle{fancy}
\fancyfoot[R]{\today}\renewcommand\headrulewidth{1pt}
\fancyhead[L]{ESTP / Bachelor 1 / \the\year}
\fancyhead[R]{Composition de math\'ematiques}
\renewcommand\footrulewidth{1pt}
%\fancyfoot[L]{Math\' ematiques : les polyn\' emes}
%\fancyfoot[R]{\thepage}


\begin{document}

\vspace{2em}

\begin{center}\bf {\Large\underline{Composition de MATH\'EMATIQUES  (3h) 25 points}}\end{center}
%\begin{center}\bf {{\small Rendre le sujet avec votre copie si le graphe de l'exercice 4 pr\' esente la courbe de la fonction $f$.}}\end{center}
\hfill\break

%%%%%%%%%%
%\includegraphics[trim = gauche bas droit haut, clip, width=3cm]
%%%%%%%%%%
\vspace{2em}
\underline {\ding{110}\, \bf{Exercice 1}}\hfill  6 points\\
\vspace{1em}
%\medskip

On considère la fonction $f$ définie par $f (x) = 2 \, \text{arcsin} x + \text{arcsin}\,(1 - 2x^2)$.
\begin{enumerate}
    \item Déterminer les domaines de définition, de continuité et de dérivabilité de $f$ .
\item Calculer la dérivée $f $' de $f$, lorsqu’elle existe.
\item En déduire une expression simplifiée de $f$.

\end{enumerate}


\hfill\break
\hrule
\hfill\break

\vspace{2em}
\underline {\ding{110}\, \bf{Exercice 2}} \hfill 4 points\\
\vspace{1em}


Donner toutes les solutions des équations différentielles suivantes : 
$$
y' -3y = xe^{-x}, \qquad \qquad  y' + y = cos(2x)\,.
$$
%{\color{blue}\underline{\bf SOLUTION :}
%
%}


\hfill\break
\hrule
\hfill\break


\vspace{2em}
\underline {\ding{110}\, \bf{Exercice 3}} \hfill 5 points\\
\vspace{1em}

\begin{enumerate}
    \item Soit $f$ la fonction définie en posant : $f (x) = \dfrac{1}{x (x^2 - 1)} .$
\begin{enumerate}
    \item Quel est l'ensemble de définition de $f$ ? Dans la suite, nous noterons $D$ cet ensemble.
\item Déterminer la décomposition en éléments simples de $f$ , c'est-à-dire déterminer trois réels $a$, $b$ et $c$ tels que :
$$\forall x \in D, \qquad f (x) = \dfrac{a}{x} + \dfrac{b}{x- 1} +\dfrac{ c}{x + 1}$$
\end{enumerate}
\item Résoudre l’équation différentielle suivante :
$(E) : x(x^2 - 1)y ' + 2y = 0$
en n’oubliant pas de préciser l’intervalle de résolution (que vous êtes libre de choisir).
\item Calculer l'intégrale suivante :
$$
\int_3^5 \frac{1}{x(x^2-1)} dx
$$

\end{enumerate}



\hfill\break
\hrule
\hfill\break

\newpage

$\,$

\vspace{2em}
\underline {\ding{110}\, \bf{Exercice 4 }}\hfill  6 points
\vspace{1em}



Calculer les intégrales suivantes, on pourra intégrer par parties, faire un changement de variable ou identifier directement une primitive. 
\begin{itemize}[itemsep=0.5em]
  \item $\displaystyle I=\int_{-\pi}^\pi (1 - x) \cos(x)dx$
  \item $\displaystyle J=\int_{1}^2 \dfrac{e^{2x}}{e^x+1}dx$ en utilisant le changement de variable $u = e^x$
  \item $\displaystyle K=\int_{-\pi}^\pi \sin(x)\cos^3(x)dx$
  \item $\displaystyle L=\int_0^1 \frac{x+1}{\sqrt{x^2+2x+3}}dx$
  \item $\displaystyle M=\int_0^1 \frac{\ln(x+1)}{x+1}dx$
  \end{itemize}




\hfill\break
\hrule
\hfill\break

\vspace{2em}
\underline {\ding{110}\, \bf{Exercice 5}} \hfill 4 points\\
\vspace{1em}
\medskip

Toutes les questions suivantes sont indépendantes entre elles.

  \begin{enumerate}
  \item
  Étudier la continuité de $g : (\mathbb R_+)^2 \longrightarrow \mathbb R$ définie par
\[
\begin{cases}
\, y^x &\text{si } y > 0 ;\\
\, 0 &\text{ si } y = 0 \text{ et } x > 0 ;\\
\,  1 &\text{ si } x = y = 0.\\
\end{cases}
\]
  \item Étudier l’éventuelle limite en $(0, 0)$ de $(x, y) \mapsto \dfrac{x^3}{x^2 + y^3}$ pour $(x, y)$ dans l'ensemble de définition.

  \item Trouver l'ensemble de définition de la fonction $f$ définie sur $\mathbb{R}^2$ par $f(x)=\sqrt{x^2+y^2-16}$

\item Calculer les dérivées partielles de $f(x,y)=x^3y+5xy+2$.

\end{enumerate}


\hfill\break
\hrule
\hfill\break



\end{document}
